\documentclass[]{article}
\usepackage{geometry,tikz,lipsum,amsmath,amssymb,soul,xcolor,cancel,esint}
\tikzset{every picture/.style={line width=0.75pt}} %set default line width to 0.75pt        
\numberwithin{equation}{section}

\NewDocumentCommand{\caja}{ m O{gray!40} O{gray!10} m }{
	\begin{figure}[h]
		\centering
		\begin{tikzpicture}
			\node[anchor=text, text width=\columnwidth-1.2cm, draw, rounded corners, line width=1pt, fill=#3, inner sep=5mm] (big) {\\#4};
			\node[draw, rounded corners, line width=.5pt, fill=#2, anchor=west, xshift=5mm] (small) at (big.north west) {#1};
		\end{tikzpicture}
	\end{figure}
}
\NewDocumentCommand{\indUpDown}{ O{\Lambda} O{\mu} O{\nu} }{
	{#1}^{#2}_{\ #3}
}
\NewDocumentCommand{\indDownUp}{ O{\Lambda} O{\mu} O{\nu} }{
	{#1}_{#2}^{\ #3}
}
\newcommand{\LambdaMuNu}{\indUpDown}
\newcommand{\0}{\textcolor{gray}{0}}
\newcommand{\rbf}{\mathbf{r}}
\newcommand{\ei}{\mathbf{e}_i}
\newcommand{\ej}{\mathbf{e}_j}
\newcommand{\lc}{\epsilon_{ijk}}
\newcommand{\gr}{\nabla}
\newcommand{\di}{\nabla\cdot}
\newcommand{\ro}{\nabla\times}
\newcommand{\la}{\nabla^2}
\newcommand{\rv}{\mathbf{r}}
\newcommand{\xv}{\mathbf{x}}
\newcommand{\yv}{\mathbf{y}}
\newcommand{\zv}{\mathbf{z}}
\newcommand{\xhat}{\hat{\xv}}
\newcommand{\yhat}{\hat{\yv}}
\newcommand{\zhat}{\hat{\zv}}
\newcommand{\rhat}{\hat{\rv}}
\newcommand{\that}{\hat{\mathbf{\theta}}}
\newcommand{\phat}{\hat{\mathbf{\varphi}}}
\newcommand{\rhhat}{\hat{\mathbf{\rho}}}
\newcommand{\nhat}{\hat{\mathbf{n}}}
\newcommand{\eo}{\epsilon_{0}}
\newcommand{\uo}{\mu_{0}}
\newcommand{\ke}{\frac{1}{4\pi\eo}}
\newcommand{\rdif}{\rv - \rv'}
\newcommand{\pd}[2]{\dfrac{\partial #1}{\partial #2}}
\newcommand{\lpd}[2]{\partial #1/\partial #2}
\newcommand{\td}[2]{\dfrac{d #1}{d #2}}
\newcommand{\ltd}[2]{d #1/d #2}
\newcommand{\n}[1]{\mathbf{#1}}
%opening
\title{Ejercicio 1b}
\author{Simon Fonseca}

\begin{document}
\maketitle
In the rest frame of one B-meson:
$$p_1 = (m_B, \vec{0}), \quad p_2 = \left(\frac{m_B}{\sqrt{1-v_{rel}^2}}, \frac{m_B v_{rel}}{\sqrt{1-v_{rel}^2}}\right)$$

The scalar product is:
$$p_1 \cdot p_2 = \frac{m_B^2}{\sqrt{1-v_{rel}^2}}$$

But this must equal the same scalar product in any frame. In the lab frame:
$$p_1 \cdot p_2 = \mathcal{E}_1 \mathcal{E}_2 - \vec{p}_1 \cdot \vec{p}_2$$

We want to find $p_1 \cdot p_2 = p_{B^+} \cdot p_{B^-}$ in the lab frame. 4-momentum is conserved:
$$p_{e^-} + p_{e^+} = p_{B^+} + p_{B^-}$$

Computing $(p_{B^+} + p_{B^-})^2$:
$$(p_{B^+} + p_{B^-})^2 = p_{B^+}^2 + p_{B^-}^2 + 2p_{B^+} \cdot p_{B^-}$$

Since $p_{B^+}^2 = m_B^2$ and $p_{B^-}^2 = m_B^2$ (mass-shell condition):

$$(p_{B^+} + p_{B^-})^2 = m_B^2 + m_B^2 + 2p_{B^+} \cdot p_{B^-} = 2m_B^2 + 2p_{B^+} \cdot p_{B^-}$$

But from conservation, $(p_{B^+} + p_{B^-})^2 = (p_{e^-} + p_{e^+})^2 = s$:

$$s = 2m_B^2 + 2p_{B^+} \cdot p_{B^-}$$

Solving for the dot product:
$$p_{B^+} \cdot p_{B^-} = \frac{s - 2m_B^2}{2}$$

For massless electrons/positrons moving in opposite directions along z:
$$p_{e^-} = (E_1, 0, 0, E_1), \quad p_{e^+} = (E_2, 0, 0, -E_2)$$

Adding them:
$$p_{e^-} + p_{e^+} = (E_1 + E_2, 0, 0, E_1 - E_2)$$

For a 4-vector $p^\mu = (E, p_x, p_y, p_z)$, the invariant square is:
$$p^2 = E^2 - p_x^2 - p_y^2 - p_z^2$$

So:
$$s = (p_{e^-} + p_{e^+})^2 = (E_1 + E_2)^2 - 0^2 - 0^2 - (E_1 - E_2)^2$$

$$s = (E_1 + E_2)^2 - (E_1 - E_2)^2$$

Expanding:
$$(E_1 + E_2)^2 = E_1^2 + 2E_1E_2 + E_2^2$$
$$(E_1 - E_2)^2 = E_1^2 - 2E_1E_2 + E_2^2$$

Subtracting:
$$s = (E_1^2 + 2E_1E_2 + E_2^2) - (E_1^2 - 2E_1E_2 + E_2^2) = 4E_1E_2$$

Finally:
$$p_{B^+} \cdot p_{B^-} = \frac{4E_1E_2 - 2m_B^2}{2} = 2E_1E_2 - m_B^2$$


Therefore:
$$\frac{m_B^2}{\sqrt{1-v_{rel}^2}} = 2E_1E_2 - m_B^2$$

Solving for $v_{rel}$:
$$v_{rel} = \sqrt{1 - \frac{m_B^4}{(2E_1E_2 - m_B^2)^2}}$$
\end{document}